\documentclass{beamer}
\usetheme{Madrid}
\usecolortheme{default}
\usepackage[utf8]{inputenc}
\usepackage{amsmath}
\usepackage{amsfonts}
\usepackage{amssymb}
\usepackage{tikz}

\title{CRISC Domain D - 100 Impossible Hard Questions}
\subtitle{Information Technology and Security - Expert Level}
\author{CRISC Mastery Series}
\date{\today}

\begin{document}
	
	\begin{frame}
		\titlepage
	\end{frame}
	
	\begin{frame}{Warning}
		\begin{alertblock}{Extreme Difficulty}
			Domain D is technical but the exam tests judgment, not trivia.\\
			Every option will look plausible.\\
			Think like a risk practitioner advising on technology, not a technician.
		\end{alertblock}
	\end{frame}
	
	\begin{frame}{Outline}
		\tableofcontents
	\end{frame}
	
	% ============================================================
	\section{1. Enterprise Architecture}
	% ============================================================
	
	\begin{frame}{Q1: EA Domains}
		\begin{block}{Question}
			An organization redesigns its business processes before redesigning software. Which EA domain drives the others?
		\end{block}
		\begin{enumerate}[(A)]
			\item Technology architecture
			\item Data architecture
			\item Business architecture
			\item Application architecture
		\end{enumerate}
	\end{frame}
	
	\begin{frame}{Q1: Explanation}
		\begin{block}{Correct Answer: C}
			\textbf{Business architecture drives the rest:}
			\begin{itemize}
				\item Business architecture defines how the enterprise operates
				\item Application and data architectures follow business needs
				\item Technology architecture supports all of them
			\end{itemize}
		\end{block}
		\begin{alertblock}{Critical Trap}
			Business architecture defines the ``why''; the others implement the ``how''.
		\end{alertblock}
	\end{frame}
	
	\begin{frame}{Q2: Technology Architecture Scope}
		\begin{block}{Question}
			Which set of components belongs to technology architecture?
		\end{block}
		\begin{enumerate}[(A)]
			\item Business processes and roles
			\item Application logic and interfaces
			\item Networks, servers, storage, and operating systems
			\item Data models and repositories
		\end{enumerate}
	\end{frame}
	
	\begin{frame}{Q2: Explanation}
		\begin{block}{Correct Answer: C}
			\textbf{Technology architecture = infrastructure:}
			\begin{itemize}
				\item Networks, servers, storage, operating systems
				\item It is the physical and logical foundation
				\item Applications and data sit above it
			\end{itemize}
		\end{block}
		\begin{alertblock}{Critical Trap}
			Technology architecture = infrastructure, not applications.
		\end{alertblock}
	\end{frame}
	
	\begin{frame}{Q3: Data Architecture Purpose}
		\begin{block}{Question}
			What is the primary purpose of data architecture in EA?
		\end{block}
		\begin{enumerate}[(A)]
			\item Defining storage vendor strategy
			\item Modeling how data is structured, integrated, and used across applications
			\item Setting firewall rules
			\item Managing employee access
		\end{enumerate}
	\end{frame}
	
	\begin{frame}{Q3: Explanation}
		\begin{block}{Correct Answer: B}
			\textbf{Data architecture models data:}
			\begin{itemize}
				\item Structure, flow, and integration
				\item Enables applications to use data consistently
				\item Supports reporting and compliance
			\end{itemize}
		\end{block}
		\begin{alertblock}{Critical Trap}
			Data architecture is about modeling, not vendors.
		\end{alertblock}
	\end{frame}
	
	\begin{frame}{Q4: EA Frameworks}
		\begin{block}{Question}
			Which framework is MOST associated with enterprise architecture?
		\end{block}
		\begin{enumerate}[(A)]
			\item COBIT
			\item TOGAF
			\item ITIL
			\item ISO 9001
		\end{enumerate}
	\end{frame}
	
	\begin{frame}{Q4: Explanation}
		\begin{block}{Correct Answer: B}
			\textbf{TOGAF is the EA standard:}
			\begin{itemize}
				\item The Open Group Architecture Framework
				\item Provides method and content for EA
				\item COBIT governs IT; ITIL governs service management
			\end{itemize}
		\end{block}
		\begin{alertblock}{Critical Trap}
			COBIT governs; TOGAF structures the architecture.
		\end{alertblock}
	\end{frame}
	
	\begin{frame}{Q5: Architecture Principle Purpose}
		\begin{block}{Question}
			Why does EA define architecture principles?
		\end{block}
		\begin{enumerate}[(A)]
			\item To slow down projects
			\item To guide consistent decisions across solutions
			\item To reduce vendor competition
			\item To avoid governance
		\end{enumerate}
	\end{frame}
	
	\begin{frame}{Q5: Explanation}
		\begin{block}{Correct Answer: B}
			\textbf{Principles guide decisions:}
			\begin{itemize}
				\item Principles guide decisions
				\item They align solutions with strategy
				\item They reduce rework and drift
			\end{itemize}
		\end{block}
		\begin{alertblock}{Critical Trap}
			Principles guide, they do not block.
		\end{alertblock}
	\end{frame}
	
	\begin{frame}{Q6: CMM Level Interpretation}
		\begin{block}{Question}
			A process is described as ``documented, standardized, and proactive.'' Which CMM level?
		\end{block}
		\begin{enumerate}[(A)]
			\item Level 1 - Initial
			\item Level 2 - Repeatable
			\item Level 3 - Defined
			\item Level 5 - Optimizing
		\end{enumerate}
	\end{frame}
	
	\begin{frame}{Q6: Explanation}
		\begin{block}{Correct Answer: C}
			\textbf{Level 3 = Defined:}
			\begin{itemize}
				\item Standardized, documented, proactive
				\item Level 2 is basic repeatable
				\item Level 5 is optimizing with metrics and continuous improvement
			\end{itemize}
		\end{block}
		\begin{alertblock}{Critical Trap}
			Defined = standardized, not yet optimized.
		\end{alertblock}
	\end{frame}
	
	\begin{frame}{Q7: CMM Level 5 Indicator}
		\begin{block}{Question}
			Which indicator BEST fits CMM Level 5?
		\end{block}
		\begin{enumerate}[(A)]
			\item Documented procedures exist
			\item Teams follow ad hoc processes
			\item Metrics drive continuous improvement
			\item Controls are undocumented
		\end{enumerate}
	\end{frame}
	
	\begin{frame}{Q7: Explanation}
		\begin{block}{Correct Answer: C}
			\textbf{Level 5 = Optimizing:}
			\begin{itemize}
				\item Metrics and continuous improvement
				\item Processes adapt and improve
				\item Level 3 is defined, Level 4 is managed
			\end{itemize}
		\end{block}
		\begin{alertblock}{Critical Trap}
			Optimizing means improving, not just documenting.
		\end{alertblock}
	\end{frame}
	
	\begin{frame}{Q8: Architecture Debt}
		\begin{block}{Question}
			What is the PRIMARY risk of accumulated architecture debt?
		\end{block}
		\begin{enumerate}[(A)]
			\item Faster innovation
			\item Increasing cost and difficulty of change
			\item Better security
			\item Easier integration
		\end{enumerate}
	\end{frame}
	
	\begin{frame}{Q8: Explanation}
		\begin{block}{Correct Answer: B}
			\textbf{Architecture debt compounds:}
			\begin{itemize}
				\item Legacy decisions limit future options
				\item Changes cost more over time
				\item Innovation slows
			\end{itemize}
		\end{block}
		\begin{alertblock}{Critical Trap}
			Architecture debt compounds silently.
		\end{alertblock}
	\end{frame}
	
	\begin{frame}{Q9: EA Governance}
		\begin{block}{Question}
			Who should own EA governance decisions?
		\end{block}
		\begin{enumerate}[(A)]
			\item A single developer
			\item An architecture review board or equivalent governance body
			\item The vendor
			\item The helpdesk
		\end{enumerate}
	\end{frame}
	
	\begin{frame}{Q9: Explanation}
		\begin{block}{Correct Answer: B}
			\textbf{EA governance needs a body:}
			\begin{itemize}
				\item Architecture review boards align decisions
				\item Cross-functional participation matters
				\item Single owners lack perspective
			\end{itemize}
		\end{block}
		\begin{alertblock}{Critical Trap}
			EA decisions need governance, not individuals.
		\end{alertblock}
	\end{frame}
	
	\begin{frame}{Q10: EA Alignment}
		\begin{block}{Question}
			Which statement BEST describes EA alignment?
		\end{block}
		\begin{enumerate}[(A)]
			\item Aligning vendors
			\item Aligning technology investments with business strategy
			\item Aligning audit plans
			\item Aligning tickets
		\end{enumerate}
	\end{frame}
	
	\begin{frame}{Q10: Explanation}
		\begin{block}{Correct Answer: B}
			\textbf{EA aligns technology to strategy:}
			\begin{itemize}
				\item Alignment is with strategy
				\item It turns business goals into technology choices
				\item It prevents orphan projects
			\end{itemize}
		\end{block}
		\begin{alertblock}{Critical Trap}
			EA aligns technology to strategy, not vice versa.
		\end{alertblock}
	\end{frame}
	
	% ============================================================
	\section{2. IT Operations and Change Management}
	% ============================================================
	
	\begin{frame}{Q11: Change Management Purpose}
		\begin{block}{Question}
			What is the PRIMARY purpose of IT change management?
		\end{block}
		\begin{enumerate}[(A)]
			\item To slow down deployments
			\item To control changes so they do not introduce unmanaged risk
			\item To avoid documentation
			\item To bypass testing
		\end{enumerate}
	\end{frame}
	
	\begin{frame}{Q11: Explanation}
		\begin{block}{Correct Answer: B}
			\textbf{Change management controls risk:}
			\begin{itemize}
				\item Change is inevitable
				\item Control reduces risk and impact
				\item Speed without control increases exposure
			\end{itemize}
		\end{block}
		\begin{alertblock}{Critical Trap}
			Change control is risk control, not bureaucracy.
		\end{alertblock}
	\end{frame}
	
	\begin{frame}{Q12: Emergency Change}
		\begin{block}{Question}
			An emergency change bypasses normal approval. What governance should still apply?
		\end{block}
		\begin{enumerate}[(A)]
			\item No governance
			\item Retrospective review, documentation, and controlled rollback readiness
			\item Ignoring the change
			\item No testing
		\end{enumerate}
	\end{frame}
	
	\begin{frame}{Q12: Explanation}
		\begin{block}{Correct Answer: B}
			\textbf{Emergency still needs controls:}
			\begin{itemize}
				\item Emergency must be documented after the fact
				\item Retrospective review catches issues
				\item Rollback readiness reduces impact
			\end{itemize}
		\end{block}
		\begin{alertblock}{Critical Trap}
			Emergency is not ``no governance''.
		\end{alertblock}
	\end{frame}
	
	\begin{frame}{Q13: CAB Composition}
		\begin{block}{Question}
			Who should sit on the Change Advisory Board?
		\end{block}
		\begin{enumerate}[(A)]
			\item Only developers
			\item Only IT operations
			\item A cross-functional group including operations, security, and business
			\item Only the vendor
		\end{enumerate}
	\end{frame}
	
	\begin{frame}{Q13: Explanation}
		\begin{block}{Correct Answer: C}
			\textbf{CAB must be cross-functional:}
			\begin{itemize}
				\item Operations evaluates impact
				\item Security evaluates risk
				\item Business evaluates priority
			\end{itemize}
		\end{block}
		\begin{alertblock}{Critical Trap}
			CAB needs cross-functional voice.
		\end{alertblock}
	\end{frame}
	
	\begin{frame}{Q14: Incident vs Problem}
		\begin{block}{Question}
			What is the MAIN distinction between incident and problem management?
		\end{block}
		\begin{enumerate}[(A)]
			\item They are identical
			\item Incidents restore service; problems find root causes
			\item Problems only track tickets
			\item Incidents only apply to networks
		\end{enumerate}
	\end{frame}
	
	\begin{frame}{Q14: Explanation}
		\begin{block}{Correct Answer: B}
			\textbf{Incident vs Problem:}
			\begin{itemize}
				\item Incidents restore service quickly
				\item Problems eliminate root causes
				\item Both are needed
			\end{itemize}
		\end{block}
		\begin{alertblock}{Critical Trap}
			Incident fixes the symptom; problem fixes the cause.
		\end{alertblock}
	\end{frame}
	
	\begin{frame}{Q15: IT Asset Management}
		\begin{block}{Question}
			What is the PRIMARY benefit of IT asset management?
		\end{block}
		\begin{enumerate}[(A)]
			\item Lower print costs
			\item Knowing what you have and where it is, so you can protect it
			\item Better lunch menu
			\item Faster emails
		\end{enumerate}
	\end{frame}
	
	\begin{frame}{Q15: Explanation}
		\begin{block}{Correct Answer: B}
			\textbf{You protect what you know:}
			\begin{itemize}
				\item Asset inventory supports risk assessment
				\item Unknown assets are unmanaged
				\item Inventory enables lifecycle management
			\end{itemize}
		\end{block}
		\begin{alertblock}{Critical Trap}
			Unknown assets are unmanaged assets.
		\end{alertblock}
	\end{frame}
	
	\begin{frame}{Q16: Configuration Management}
		\begin{block}{Question}
			Why is configuration management essential to security?
		\end{block}
		\begin{enumerate}[(A)]
			\item It speeds up printing
			\item It provides known, secure baselines and detects drift
			\item It reduces HR work
			\item It avoids audits
		\end{enumerate}
	\end{frame}
	
	\begin{frame}{Q16: Explanation}
		\begin{block}{Correct Answer: B}
			\textbf{Baselines enable security:}
			\begin{itemize}
				\item Baselines define secure states
				\item Drift detection catches misconfigurations
				\item Standardization reduces risk
			\end{itemize}
		\end{block}
		\begin{alertblock}{Critical Trap}
			Configurations drift silently; baselines detect them.
		\end{alertblock}
	\end{frame}
	
	\begin{frame}{Q17: Release Management}
		\begin{block}{Question}
			What is the MAIN goal of release management?
		\end{block}
		\begin{enumerate}[(A)]
			\item Rolling out changes without testing
			\item Coordinating packaging and deployment to minimize risk and disruption
			\item Avoiding business involvement
			\item Reducing budgets
		\end{enumerate}
	\end{frame}
	
	\begin{frame}{Q17: Explanation}
		\begin{block}{Correct Answer: B}
			\textbf{Release management coordinates deployment:}
			\begin{itemize}
				\item Packaging, testing, deployment coordination
				\item Minimizes risk and disruption
				\item Business involvement is essential
			\end{itemize}
		\end{block}
		\begin{alertblock}{Critical Trap}
			Release management is coordination, not code.
		\end{alertblock}
	\end{frame}
	
	\begin{frame}{Q18: Parallel Run}
		\begin{block}{Question}
			When is a parallel run MOST appropriate?
		\end{block}
		\begin{enumerate}[(A)]
			\item For minor cosmetic changes
			\item For critical systems where downtime would cause severe impact
			\item When budget is unlimited
			\item When no testing is planned
		\end{enumerate}
	\end{frame}
	
	\begin{frame}{Q18: Explanation}
		\begin{block}{Correct Answer: B}
			\textbf{Parallel runs protect critical systems:}
			\begin{itemize}
				\item Both systems run simultaneously
				\item Highest assurance but highest cost
				\item Suitable for critical transitions
			\end{itemize}
		\end{block}
		\begin{alertblock}{Critical Trap}
			Parallel runs are expensive but safest.
		\end{alertblock}
	\end{frame}
	
	\begin{frame}{Q19: Phased Rollout}
		\begin{block}{Question}
			Which changeover method BEST suits independent modules with no cross-dependencies?
		\end{block}
		\begin{enumerate}[(A)]
			\item Abrupt
			\item Parallel
			\item Phased
			\item Big bang
		\end{enumerate}
	\end{frame}
	
	\begin{frame}{Q19: Explanation}
		\begin{block}{Correct Answer: C}
			\textbf{Phased suits independent modules:}
			\begin{itemize}
				\item Modules deployed one at a time
				\item Reduced risk of full outage
				\item Independent modules make phasing feasible
			\end{itemize}
		\end{block}
		\begin{alertblock}{Critical Trap}
			Independent modules means phased rollout.
		\end{alertblock}
	\end{frame}
	
	\begin{frame}{Q20: Abrupt Changeover}
		\begin{block}{Question}
			Which situation MOST justifies an abrupt changeover?
		\end{block}
		\begin{enumerate}[(A)]
			\item Critical 24/7 system
			\item Minor, easily reversible change with low impact
			\item High-risk financial trading system
			\item Regulatory reporting system
		\end{enumerate}
	\end{frame}
	
	\begin{frame}{Q20: Explanation}
		\begin{block}{Correct Answer: B}
			\textbf{Abrupt is for low-risk changes:}
			\begin{itemize}
				\item Abrupt is highest risk
				\item Suitable only for low-impact, reversible changes
				\item Critical systems need safer changeover
			\end{itemize}
		\end{block}
		\begin{alertblock}{Critical Trap}
			Abrupt is highest risk. Only for low-risk changes.
		\end{alertblock}
	\end{frame}
	
	\begin{frame}{Q21: Rollback Plan}
		\begin{block}{Question}
			Why must a rollback plan exist before deployment?
		\end{block}
		\begin{enumerate}[(A)]
			\item To slow down the project
			\item To provide a safe recovery path if the change fails
			\item To satisfy vendors
			\item To avoid testing
		\end{enumerate}
	\end{frame}
	
	\begin{frame}{Q21: Explanation}
		\begin{block}{Correct Answer: B}
			\textbf{Rollback protects operations:}
			\begin{itemize}
				\item Changes can fail
				\item Rollback restores service quickly
				\item Planning prevents panic
			\end{itemize}
		\end{block}
		\begin{alertblock}{Critical Trap}
			No rollback means risky deployment.
		\end{alertblock}
	\end{frame}
	
	\begin{frame}{Q22: Post-Implementation Review}
		\begin{block}{Question}
			Which outcome signals a successful post-implementation review?
		\end{block}
		\begin{enumerate}[(A)]
			\item No documentation
			\item Documented lessons and improved practices for future changes
			\item Blame assigned
			\item Budget increased
		\end{enumerate}
	\end{frame}
	
	\begin{frame}{Q22: Explanation}
		\begin{block}{Correct Answer: B}
			\textbf{PIR drives improvement:}
			\begin{itemize}
				\item Lessons documented and used
				\item Feedback improves future changes
				\item Blame discourages learning
			\end{itemize}
		\end{block}
		\begin{alertblock}{Critical Trap}
			PIR without lessons is paperwork.
		\end{alertblock}
	\end{frame}
	
	\begin{frame}{Q23: Production Data in Tests}
		\begin{block}{Question}
			What is the MAIN risk of using production data for testing?
		\end{block}
		\begin{enumerate}[(A)]
			\item Faster tests
			\item Unauthorized exposure of sensitive data
			\item Reduced cost
			\item Better collaboration
		\end{enumerate}
	\end{frame}
	
	\begin{frame}{Q23: Explanation}
		\begin{block}{Correct Answer: B}
			\textbf{Sensitive data exposure:}
			\begin{itemize}
				\item Production data contains PII, PHI, and secrets
				\item Test environments are less secure
				\item Masking or synthetic data avoids risk
			\end{itemize}
		\end{block}
		\begin{alertblock}{Critical Trap}
			Production data in tests leaks data.
		\end{alertblock}
	\end{frame}
	
	\begin{frame}{Q24: Change Ticket Traceability}
		\begin{block}{Question}
			Why is change ticket traceability important for audits?
		\end{block}
		\begin{enumerate}[(A)]
			\item It saves money
			\item It provides evidence of approval, testing, and rollback
			\item It looks professional
			\item It avoids testing
		\end{enumerate}
	\end{frame}
	
	\begin{frame}{Q24: Explanation}
		\begin{block}{Correct Answer: B}
			\textbf{Traceability shows governance:}
			\begin{itemize}
				\item Approval, testing, rollback evidence
				\item Supports audits and incident review
				\item Shows due care
			\end{itemize}
		\end{block}
		\begin{alertblock}{Critical Trap}
			No trace means no evidence of control.
		\end{alertblock}
	\end{frame}
	
	\begin{frame}{Q25: Shift-Left Security}
		\begin{block}{Question}
			What does ``shift-left security'' mean in IT operations?
		\end{block}
		\begin{enumerate}[(A)]
			\item Security moves to the left side of the building
			\item Security is integrated earlier in the development and change process
			\item Security is postponed
			\item Security is outsourced
		\end{enumerate}
	\end{frame}
	
	\begin{frame}{Q25: Explanation}
		\begin{block}{Correct Answer: B}
			\textbf{Shift-left = earlier security:}
			\begin{itemize}
				\item Security in design and build
				\item Cheaper to fix early
				\item Reduces downstream risk
			\end{itemize}
		\end{block}
		\begin{alertblock}{Critical Trap}
			Shift-left is about timing, not location.
		\end{alertblock}
	\end{frame}
	
	% ============================================================
	\section{3. Enterprise Resiliency}
	% ============================================================
	
	\begin{frame}{Q26: Resiliency vs Recovery}
		\begin{block}{Question}
			Which BEST distinguishes resiliency from recovery?
		\end{block}
		\begin{enumerate}[(A)]
			\item They are the same thing
			\item Resiliency avoids or limits disruption; recovery restores after disruption
			\item Recovery avoids disruption; resiliency restores
			\item Resiliency only applies to hardware
		\end{enumerate}
	\end{frame}
	
	\begin{frame}{Q26: Explanation}
		\begin{block}{Correct Answer: B}
			\textbf{Resiliency before, recovery after:}
			\begin{itemize}
				\item Resiliency prevents or limits disruption
				\item Recovery restores after an event
				\item Both are needed
			\end{itemize}
		\end{block}
		\begin{alertblock}{Critical Trap}
			Resiliency anticipates; recovery restores.
		\end{alertblock}
	\end{frame}
	
	\begin{frame}{Q27: Business Continuity vs Disaster Recovery}
		\begin{block}{Question}
			Which statement is TRUE?
		\end{block}
		\begin{enumerate}[(A)]
			\item BC and DR are identical
			\item DR is a subset of BC
			\item BC is a subset of DR
			\item Both only apply to networks
		\end{enumerate}
	\end{frame}
	
	\begin{frame}{Q27: Explanation}
		\begin{block}{Correct Answer: B}
			\textbf{DR is a subset of BC:}
			\begin{itemize}
				\item BC covers processes, people, and facilities
				\item DR restores IT systems and data
				\item DR supports broader BC objectives
			\end{itemize}
		\end{block}
		\begin{alertblock}{Critical Trap}
			BC is broader; DR is one component.
		\end{alertblock}
	\end{frame}
	
	\begin{frame}{Q28: RTO vs MTD}
		\begin{block}{Question}
			An application must be restored within 2 hours, but the business could survive 12 hours of outage. Which statement is correct?
		\end{block}
		\begin{enumerate}[(A)]
			\item RTO = 12h, MTD = 2h
			\item RTO = 2h, MTD = 12h
			\item RTO = MTD
			\item RTO and MTD are the same
		\end{enumerate}
	\end{frame}
	
	\begin{frame}{Q28: Explanation}
		\begin{block}{Correct Answer: B}
			\textbf{RTO is inside MTD:}
			\begin{itemize}
				\item RTO is the recovery target (2h)
				\item MTD is the maximum tolerable outage (12h)
				\item RTO should be less than or equal to MTD
			\end{itemize}
		\end{block}
		\begin{alertblock}{Critical Trap}
			RTO must be less than or equal to MTD.
		\end{alertblock}
	\end{frame}
	
	\begin{frame}{Q29: RPO Example}
		\begin{block}{Question}
			RPO is 30 minutes. Backups run hourly. What is the issue?
		\end{block}
		\begin{enumerate}[(A)]
			\item No issue
			\item Backups may lose up to 60 minutes of data, exceeding the RPO
			\item RPO is too short
			\item Backups run too frequently
		\end{enumerate}
	\end{frame}
	
	\begin{frame}{Q29: Explanation}
		\begin{block}{Correct Answer: B}
			\textbf{Backup frequency must meet RPO:}
			\begin{itemize}
				\item Hourly backups = up to 60 minutes data loss
				\item RPO is 30 minutes
				\item Need more frequent backups or replication
			\end{itemize}
		\end{block}
		\begin{alertblock}{Critical Trap}
			Backup cadence must be at least as frequent as RPO.
		\end{alertblock}
	\end{frame}
	
	\begin{frame}{Q30: MTD Significance}
		\begin{block}{Question}
			Why is MTD a critical metric?
		\end{block}
		\begin{enumerate}[(A)]
			\item It sets the recovery target
			\item It defines the maximum tolerable outage before severe impact
			\item It measures data loss
			\item It measures backup speed
		\end{enumerate}
	\end{frame}
	
	\begin{frame}{Q30: Explanation}
		\begin{block}{Correct Answer: B}
			\textbf{MTD defines severe impact:}
			\begin{itemize}
				\item Beyond MTD, impact becomes unacceptable
				\item MTD is the ceiling
				\item RTO is the target inside MTD
			\end{itemize}
		\end{block}
		\begin{alertblock}{Critical Trap}
			MTD is the ceiling; RTO is inside.
		\end{alertblock}
	\end{frame}
	
	\begin{frame}{Q31: Failover vs Failback}
		\begin{block}{Question}
			What is failback?
		\end{block}
		\begin{enumerate}[(A)]
			\item Switching to the secondary site
			\item Returning operations to the primary site after recovery
			\item Reverting a change
			\item Rolling back a backup
		\end{enumerate}
	\end{frame}
	
	\begin{frame}{Q31: Explanation}
		\begin{block}{Correct Answer: B}
			\textbf{Failback returns to primary:}
			\begin{itemize}
				\item Failover = move to secondary
				\item Failback = return to primary
				\item Both must be tested
			\end{itemize}
		\end{block}
		\begin{alertblock}{Critical Trap}
			Failover moves; failback returns.
		\end{alertblock}
	\end{frame}
	
	\begin{frame}{Q32: DR Test Types}
		\begin{block}{Question}
			Which DR test provides the HIGHEST confidence?
		\end{block}
		\begin{enumerate}[(A)]
			\item Review of the DR plan
			\item Tabletop exercise
			\item Full failover to alternate site with real services
			\item Walkthrough
		\end{enumerate}
	\end{frame}
	
	\begin{frame}{Q32: Explanation}
		\begin{block}{Correct Answer: C}
			\textbf{Full failover = highest confidence:}
			\begin{itemize}
				\item Tests real recovery
				\item Requires more resources
				\item Tabletop and walkthrough are cheaper but lower assurance
			\end{itemize}
		\end{block}
		\begin{alertblock}{Critical Trap}
			Full failover = highest assurance, highest cost.
		\end{alertblock}
	\end{frame}
	
	\begin{frame}{Q33: Backup 3-2-1 Rule}
		\begin{block}{Question}
			Which is the 3-2-1 backup rule?
		\end{block}
		\begin{enumerate}[(A)]
			\item 3 admins, 2 vendors, 1 server
			\item 3 copies of data, 2 different media, 1 offsite
			\item 3 backups per day
			\item 3 sites, 2 networks, 1 admin
		\end{enumerate}
	\end{frame}
	
	\begin{frame}{Q33: Explanation}
		\begin{block}{Correct Answer: B}
			\textbf{3-2-1 = copies, media, offsite:}
			\begin{itemize}
				\item 3 copies of data
				\item 2 different media types
				\item 1 offsite copy
			\end{itemize}
		\end{block}
		\begin{alertblock}{Critical Trap}
			3-2-1 = copies, media, offsite.
		\end{alertblock}
	\end{frame}
	
	\begin{frame}{Q34: Air Gap Backup}
		\begin{block}{Question}
			Why are air-gapped backups important against ransomware?
		\end{block}
		\begin{enumerate}[(A)]
			\item They are cheaper
			\item They isolate backups from network attacks and encryption
			\item They are faster
			\item They do not need testing
		\end{enumerate}
	\end{frame}
	
	\begin{frame}{Q34: Explanation}
		\begin{block}{Correct Answer: B}
			\textbf{Air gap blocks ransomware:}
			\begin{itemize}
				\item Ransomware encrypts connected backups
				\item Air-gapped backups stay clean
				\item Crucial for recovery
			\end{itemize}
		\end{block}
		\begin{alertblock}{Critical Trap}
			Ransomware can encrypt connected backups.
		\end{alertblock}
	\end{frame}
	
	\begin{frame}{Q35: Recovery Testing}
		\begin{block}{Question}
			Backups exist but have never been restored. What is the KEY risk?
		\end{block}
		\begin{enumerate}[(A)]
			\item Slower backups
			\item Restore may fail when needed
			\item Higher storage costs
			\item Less data
		\end{enumerate}
	\end{frame}
	
	\begin{frame}{Q35: Explanation}
		\begin{block}{Correct Answer: B}
			\textbf{Untested restore = unreliable backup:}
			\begin{itemize}
				\item Backups can be corrupt or incomplete
				\item Restore tests reveal issues
				\item Restore capability matters, not backup count
			\end{itemize}
		\end{block}
		\begin{alertblock}{Critical Trap}
			Backups are useless if restore fails.
		\end{alertblock}
	\end{frame}
	
	% ============================================================
	\section{4. Data Lifecycle Management}
	% ============================================================
	
	\begin{frame}{Q36: Data Classification Purpose}
		\begin{block}{Question}
			What is the PRIMARY purpose of data classification?
		\end{block}
		\begin{enumerate}[(A)]
			\item Reduce storage
			\item Determine the level of protection required
			\item Simplify backups
			\item Speed up reports
		\end{enumerate}
	\end{frame}
	
	\begin{frame}{Q36: Explanation}
		\begin{block}{Correct Answer: B}
			\textbf{Classification drives protection:}
			\begin{itemize}
				\item Sensitivity determines controls
				\item Enables proportionate security
				\item Higher class means stronger controls
			\end{itemize}
		\end{block}
		\begin{alertblock}{Critical Trap}
			Classification determines controls, not storage.
		\end{alertblock}
	\end{frame}
	
	\begin{frame}{Q37: Data Labeling}
		\begin{block}{Question}
			Why is data labeling important after classification?
		\end{block}
		\begin{enumerate}[(A)]
			\item Aesthetics
			\item It makes classification actionable at the data level
			\item It speeds up backups
			\item It reduces costs
		\end{enumerate}
	\end{frame}
	
	\begin{frame}{Q37: Explanation}
		\begin{block}{Correct Answer: B}
			\textbf{Labeling operationalizes classification:}
			\begin{itemize}
				\item Labels identify sensitivity at rest and in transit
				\item DLP tools can enforce based on labels
				\item Unlabeled data bypasses controls
			\end{itemize}
		\end{block}
		\begin{alertblock}{Critical Trap}
			Classification without labeling is incomplete.
		\end{alertblock}
	\end{frame}
	
	\begin{frame}{Q38: Data Lifecycle Stages}
		\begin{block}{Question}
			Which stage handles disposal of data?
		\end{block}
		\begin{enumerate}[(A)]
			\item Creation
			\item Use
			\item Archiving
			\item Destruction
		\end{enumerate}
	\end{frame}
	
	\begin{frame}{Q38: Explanation}
		\begin{block}{Correct Answer: D}
			\textbf{Destruction is disposal:}
			\begin{itemize}
				\item Secure disposal of data
				\item Meets legal and regulatory requirements
				\item Archiving is retention, not destruction
			\end{itemize}
		\end{block}
		\begin{alertblock}{Critical Trap}
			Destruction is the disposal stage.
		\end{alertblock}
	\end{frame}
	
	\begin{frame}{Q39: Data Retention Policy}
		\begin{block}{Question}
			What does a data retention policy specify?
		\end{block}
		\begin{enumerate}[(A)]
			\item Data value only
			\item How long data must be retained and when destroyed
			\item Vendor preferences
			\item Analytics needs
		\end{enumerate}
	\end{frame}
	
	\begin{frame}{Q39: Explanation}
		\begin{block}{Correct Answer: B}
			\textbf{Retention defines duration and destruction:}
			\begin{itemize}
				\item Retain for legal, regulatory, business needs
				\item Destroy when no longer needed
				\item Reduce exposure and cost
			\end{itemize}
		\end{block}
		\begin{alertblock}{Critical Trap}
			Retention is how long, not value.
		\end{alertblock}
	\end{frame}
	
	\begin{frame}{Q40: Data Remanence}
		\begin{block}{Question}
			What is data remanence and why does it matter?
		\end{block}
		\begin{enumerate}[(A)]
			\item Leftover data after deletion that can be recovered
			\item Data stored in RAM
			\item Data in backups
			\item Data in logs
		\end{enumerate}
	\end{frame}
	
	\begin{frame}{Q40: Explanation}
		\begin{block}{Correct Answer: A}
			\textbf{Deletion is not destruction:}
			\begin{itemize}
				\item Residual data can be recovered
				\item Proper destruction uses standards
				\item Sanitization or physical destruction
			\end{itemize}
		\end{block}
		\begin{alertblock}{Critical Trap}
			Deleting files does not remove data.
		\end{alertblock}
	\end{frame}
	
	\begin{frame}{Q41: DLP Purpose}
		\begin{block}{Question}
			What is the PRIMARY purpose of a Data Loss Prevention program?
		\end{block}
		\begin{enumerate}[(A)]
			\item Prevent unauthorized data exfiltration
			\item Slow down emails
			\item Store data cheaply
			\item Avoid backups
		\end{enumerate}
	\end{frame}
	
	\begin{frame}{Q41: Explanation}
		\begin{block}{Correct Answer: A}
			\textbf{DLP prevents exfiltration:}
			\begin{itemize}
				\item Detects and blocks sensitive data leaving
				\item Uses classification and content inspection
				\item Reduces insider and exfiltration risk
			\end{itemize}
		\end{block}
		\begin{alertblock}{Critical Trap}
			DLP prevents loss, not storage.
		\end{alertblock}
	\end{frame}
	
	\begin{frame}{Q42: Data Minimization}
		\begin{block}{Question}
			What does data minimization require?
		\end{block}
		\begin{enumerate}[(A)]
			\item Collect as much as possible
			\item Collect only the data needed for the stated purpose
			\item Keep data forever
			\item Store data without encryption
		\end{enumerate}
	\end{frame}
	
	\begin{frame}{Q42: Explanation}
		\begin{block}{Correct Answer: B}
			\textbf{Minimization reduces exposure:}
			\begin{itemize}
				\item Less data collected means less risk
				\item Privacy regulations often require it
				\item Simplifies compliance
			\end{itemize}
		\end{block}
		\begin{alertblock}{Critical Trap}
			Data minimization means less is safer.
		\end{alertblock}
	\end{frame}
	
	\begin{frame}{Q43: Data Governance vs Management}
		\begin{block}{Question}
			Which statement is TRUE about data governance vs data management?
		\end{block}
		\begin{enumerate}[(A)]
			\item They are identical
			\item Governance sets policy; management executes
			\item Management sets policy; governance executes
			\item Both apply only to backups
		\end{enumerate}
	\end{frame}
	
	\begin{frame}{Q43: Explanation}
		\begin{block}{Correct Answer: B}
			\textbf{Governance vs management:}
			\begin{itemize}
				\item Governance sets direction and accountability
				\item Management executes and enforces
				\item Roles must be clear
			\end{itemize}
		\end{block}
		\begin{alertblock}{Critical Trap}
			Governance sets policy; management enforces.
		\end{alertblock}
	\end{frame}
	
	\begin{frame}{Q44: Archiving Purpose}
		\begin{block}{Question}
			Why is data archived rather than deleted immediately?
		\end{block}
		\begin{enumerate}[(A)]
			\item To save storage
			\item To meet legal, regulatory, or contractual retention requirements
			\item To avoid audits
			\item To slow down systems
		\end{enumerate}
	\end{frame}
	
	\begin{frame}{Q44: Explanation}
		\begin{block}{Correct Answer: B}
			\textbf{Archiving meets retention needs:}
			\begin{itemize}
				\item Legal and contractual requirements
				\item Preserve when not in active use
				\item Easier and cheaper than active storage
			\end{itemize}
		\end{block}
		\begin{alertblock}{Critical Trap}
			Archiving is about retention, not active use.
		\end{alertblock}
	\end{frame}
	
	\begin{frame}{Q45: Secure Destruction}
		\begin{block}{Question}
			Which is MOST appropriate for destroying data on decommissioned drives?
		\end{block}
		\begin{enumerate}[(A)]
			\item Deleting files
			\item Formatting
			\item Standards-based sanitization or physical destruction
			\item Reusing drives
		\end{enumerate}
	\end{frame}
	
	\begin{frame}{Q45: Explanation}
		\begin{block}{Correct Answer: C}
			\textbf{Standards-based destruction:}
			\begin{itemize}
				\item Use recognized standards such as NIST 800-88
				\item Deleting or formatting leaves data recoverable
				\item Physical destruction may be required for high sensitivity
			\end{itemize}
		\end{block}
		\begin{alertblock}{Critical Trap}
			Deleting is not destroying.
		\end{alertblock}
	\end{frame}
	
	% ============================================================
	\section{5. SDLC and Emerging Technologies}
	% ============================================================
	
	\begin{frame}{Q46: SDLC Phases}
		\begin{block}{Question}
			Which phase of the SDLC includes budgeting and requirements definition?
		\end{block}
		\begin{enumerate}[(A)]
			\item Development
			\item Implementation
			\item Initiation
			\item Disposal
		\end{enumerate}
	\end{frame}
	
	\begin{frame}{Q46: Explanation}
		\begin{block}{Correct Answer: C}
			\textbf{Initiation sets the foundation:}
			\begin{itemize}
				\item Feasibility and scope
				\item Requirements and budget
				\item Everything else depends on it
			\end{itemize}
		\end{block}
		\begin{alertblock}{Critical Trap}
			Initiation includes budgeting and requirements.
		\end{alertblock}
	\end{frame}
	
	\begin{frame}{Q47: SDLC Disposal}
		\begin{block}{Question}
			What is the primary focus of the SDLC disposal phase?
		\end{block}
		\begin{enumerate}[(A)]
			\item Coding new features
			\item Secure destruction of data and decommissioning
			\item Training users
			\item Writing requirements
		\end{enumerate}
	\end{frame}
	
	\begin{frame}{Q47: Explanation}
		\begin{block}{Correct Answer: B}
			\textbf{Disposal is secure end-of-life:}
			\begin{itemize}
				\item Destroy data securely
				\item Decommission systems
				\item Preserve records as required
			\end{itemize}
		\end{block}
		\begin{alertblock}{Critical Trap}
			Disposal is not just deleting the system.
		\end{alertblock}
	\end{frame}
	
	\begin{frame}{Q48: SDLC User Training}
		\begin{block}{Question}
			Which SDLC phase typically includes user training?
		\end{block}
		\begin{enumerate}[(A)]
			\item Initiation
			\item Development
			\item Implementation
			\item Disposal
		\end{enumerate}
	\end{frame}
	
	\begin{frame}{Q48: Explanation}
		\begin{block}{Correct Answer: C}
			\textbf{Training happens at implementation:}
			\begin{itemize}
				\item Users trained before or during rollout
				\item Implementation includes adoption
				\item Development builds, implementation deploys
			\end{itemize}
		\end{block}
		\begin{alertblock}{Critical Trap}
			User training is part of implementation.
		\end{alertblock}
	\end{frame}
	
	\begin{frame}{Q49: Project Risk vs SDLC Risk}
		\begin{block}{Question}
			Which statement BEST distinguishes project risk from SDLC risk?
		\end{block}
		\begin{enumerate}[(A)]
			\item They are identical
			\item Project risk focuses on objectives; SDLC risk focuses on the development process
			\item SDLC risk focuses on budgets only
			\item Project risk applies only to IT
		\end{enumerate}
	\end{frame}
	
	\begin{frame}{Q49: Explanation}
		\begin{block}{Correct Answer: B}
			\textbf{Two kinds of risk:}
			\begin{itemize}
				\item Project risk = achieving objectives
				\item SDLC risk = defects in design, coding, testing
				\item Both must be managed
			\end{itemize}
		\end{block}
		\begin{alertblock}{Critical Trap}
			Project risk is about the project; SDLC risk is about the product.
		\end{alertblock}
	\end{frame}
	
	\begin{frame}{Q50: Accreditation vs Certification}
		\begin{block}{Question}
			Which statement is TRUE about accreditation and certification?
		\end{block}
		\begin{enumerate}[(A)]
			\item They are the same
			\item Accreditation can be internal or external; certification is external
			\item Only accreditation exists
			\item Only certification is required by law
		\end{enumerate}
	\end{frame}
	
	\begin{frame}{Q50: Explanation}
		\begin{block}{Correct Answer: B}
			\textbf{Accreditation vs certification:}
			\begin{itemize}
				\item Accreditation = evaluation of compliance to standards
				\item Certification = formal recognition by external body
				\item Both provide assurance
			\end{itemize}
		\end{block}
		\begin{alertblock}{Critical Trap}
			Accreditation can be internal; certification is external.
		\end{alertblock}
	\end{frame}
	
	\begin{frame}{Q51: BYOD Risk}
		\begin{block}{Question}
			What is the PRIMARY security risk of BYOD?
		\end{block}
		\begin{enumerate}[(A)]
			\item Cheaper devices
			\item Loss of control over devices containing organizational data
			\item Higher morale
			\item Faster approvals
		\end{enumerate}
	\end{frame}
	
	\begin{frame}{Q51: Explanation}
		\begin{block}{Correct Answer: B}
			\textbf{BYOD reduces control:}
			\begin{itemize}
				\item Personal devices are outside enterprise control
				\item Data leakage and unpatched devices
				\item MDM and policy are essential
			\end{itemize}
		\end{block}
		\begin{alertblock}{Critical Trap}
			BYOD trades control for convenience.
		\end{alertblock}
	\end{frame}
	
	\begin{frame}{Q52: IoT Risk}
		\begin{block}{Question}
			Which is the MOST significant security concern for IoT devices?
		\end{block}
		\begin{enumerate}[(A)]
			\item They are expensive
			\item Weak default security and limited patching
			\item They consume too much power
			\item They are new
		\end{enumerate}
	\end{frame}
	
	\begin{frame}{Q52: Explanation}
		\begin{block}{Correct Answer: B}
			\textbf{Weak defaults and limited patching:}
			\begin{itemize}
				\item IoT devices often lack strong security
				\item Patch support is often minimal
				\item Network segmentation helps
			\end{itemize}
		\end{block}
		\begin{alertblock}{Critical Trap}
			IoT is risky because it is unmanaged.
		\end{alertblock}
	\end{frame}
	
	\begin{frame}{Q53: AI Risk}
		\begin{block}{Question}
			Which AI risk MOST affects decision-making reliability?
		\end{block}
		\begin{enumerate}[(A)]
			\item Cost of GPUs
			\item Bias and lack of explainability
			\item Number of users
			\item Hardware age
		\end{enumerate}
	\end{frame}
	
	\begin{frame}{Q53: Explanation}
		\begin{block}{Correct Answer: B}
			\textbf{Bias and opacity:}
			\begin{itemize}
				\item Biased data produces biased outcomes
				\item Unexplainable models cannot be governed
				\item Explainability is often required
			\end{itemize}
		\end{block}
		\begin{alertblock}{Critical Trap}
			AI risk is often bias and opacity.
		\end{alertblock}
	\end{frame}
	
	\begin{frame}{Q54: Blockchain Risk}
		\begin{block}{Question}
			Which is a PRIMARY risk of blockchain-based systems?
		\end{block}
		\begin{enumerate}[(A)]
			\item Immutability means errors cannot be easily corrected
			\item It uses too much paper
			\item It is too new
			\item It requires no governance
		\end{enumerate}
	\end{frame}
	
	\begin{frame}{Q54: Explanation}
		\begin{block}{Correct Answer: A}
			\textbf{Immutability has trade-offs:}
			\begin{itemize}
				\item Records cannot be changed easily
				\item Bad data persists
				\item Smart contract bugs are hard to fix
			\end{itemize}
		\end{block}
		\begin{alertblock}{Critical Trap}
			Immutability is both strength and risk.
		\end{alertblock}
	\end{frame}
	
	\begin{frame}{Q55: Quantum Risk}
		\begin{block}{Question}
			What is the MAIN risk quantum computing poses to security?
		\end{block}
		\begin{enumerate}[(A)]
			\item Faster browsing
			\item Breaking current public-key cryptography
			\item Cheaper storage
			\item Faster backups
		\end{enumerate}
	\end{frame}
	
	\begin{frame}{Q55: Explanation}
		\begin{block}{Correct Answer: B}
			\textbf{Quantum breaks public-key crypto:}
			\begin{itemize}
				\item RSA and ECC may be at risk
				\item Post-quantum cryptography is emerging
				\item Harvesting encrypted data today may reveal it later
			\end{itemize}
		\end{block}
		\begin{alertblock}{Critical Trap}
			Quantum threatens current encryption.
		\end{alertblock}
	\end{frame}
	
	\begin{frame}{Q56: Secure SDLC}
		\begin{block}{Question}
			What does a secure SDLC include?
		\end{block}
		\begin{enumerate}[(A)]
			\item Security only at the end
			\item Security activities embedded at each phase
			\item Only testing
			\item Only documentation
		\end{enumerate}
	\end{frame}
	
	\begin{frame}{Q56: Explanation}
		\begin{block}{Correct Answer: B}
			\textbf{Security throughout the lifecycle:}
			\begin{itemize}
				\item Requirements, design, build, test, deploy, operate
				\item Early security is cheaper
				\item End-only security misses risks
			\end{itemize}
		\end{block}
		\begin{alertblock}{Critical Trap}
			Secure SDLC is security in every phase.
		\end{alertblock}
	\end{frame}
	
	\begin{frame}{Q57: Threat Modeling Timing}
		\begin{block}{Question}
			When should threat modeling be performed?
		\end{block}
		\begin{enumerate}[(A)]
			\item Only after go-live
			\item During design, before implementation
			\item Only during testing
			\item Only after an incident
		\end{enumerate}
	\end{frame}
	
	\begin{frame}{Q57: Explanation}
		\begin{block}{Correct Answer: B}
			\textbf{Threat modeling in design:}
			\begin{itemize}
				\item Earliest and cheapest fix
				\item Shapes architecture and requirements
				\item Post-go-live is expensive
			\end{itemize}
		\end{block}
		\begin{alertblock}{Critical Trap}
			Threat modeling is a design activity.
		\end{alertblock}
	\end{frame}
	
	\begin{frame}{Q58: Code Review Purpose}
		\begin{block}{Question}
			What is the MAIN purpose of code review?
		\end{block}
		\begin{enumerate}[(A)]
			\item To meet deadlines
			\item To find defects and security issues before release
			\item To train juniors only
			\item To reduce code size
		\end{enumerate}
	\end{frame}
	
	\begin{frame}{Q58: Explanation}
		\begin{block}{Correct Answer: B}
			\textbf{Catch issues early:}
			\begin{itemize}
				\item Defects are cheaper to fix early
				\item Manual review complements tools
				\item Peer review improves quality
			\end{itemize}
		\end{block}
		\begin{alertblock}{Critical Trap}
			Code review is about finding defects, not metrics.
		\end{alertblock}
	\end{frame}
	
	\begin{frame}{Q59: SAST vs DAST}
		\begin{block}{Question}
			What is the MAIN difference between SAST and DAST?
		\end{block}
		\begin{enumerate}[(A)]
			\item SAST analyzes source code statically; DAST tests running applications dynamically
			\item Both are the same
			\item SAST needs no source
			\item DAST only scans binaries
		\end{enumerate}
	\end{frame}
	
	\begin{frame}{Q59: Explanation}
		\begin{block}{Correct Answer: A}
			\textbf{SAST is static; DAST is dynamic:}
			\begin{itemize}
				\item SAST scans source code without running
				\item DAST tests live running apps
				\item Both are useful and complementary
			\end{itemize}
		\end{block}
		\begin{alertblock}{Critical Trap}
			SAST analyzes code; DAST exercises runtime.
		\end{alertblock}
	\end{frame}
	
	\begin{frame}{Q60: Penetration Testing Scope}
		\begin{block}{Question}
			What MUST be defined before a penetration test begins?
		\end{block}
		\begin{enumerate}[(A)]
			\item Costs only
			\item Scope, rules of engagement, and authorization
			\item Vendor names
			\item Marketing message
		\end{enumerate}
	\end{frame}
	
	\begin{frame}{Q60: Explanation}
		\begin{block}{Correct Answer: B}
			\textbf{Scope protects everyone:}
			\begin{itemize}
				\item Scope protects boundaries
				\item Authorization protects legality
				\item Rules of engagement protect operations
			\end{itemize}
		\end{block}
		\begin{alertblock}{Critical Trap}
			Pen tests without scope are dangerous and illegal.
		\end{alertblock}
	\end{frame}
	
	% ============================================================
	\section{6. Information Security Principles}
	% ============================================================
	
	\begin{frame}{Q61: CIA Triad}
		\begin{block}{Question}
			Which of the following BEST describes the CIA triad?
		\end{block}
		\begin{enumerate}[(A)]
			\item Confidentiality, Integrity, Availability
			\item Compliance, Integrity, Access
			\item Control, Isolation, Assurance
			\item Cybersecurity, IT, Auditing
		\end{enumerate}
	\end{frame}
	
	\begin{frame}{Q61: Explanation}
		\begin{block}{Correct Answer: A}
			\textbf{CIA = Confidentiality, Integrity, Availability:}
			\begin{itemize}
				\item Confidentiality = only authorized access
				\item Integrity = accuracy and completeness
				\item Availability = access when needed
			\end{itemize}
		\end{block}
		\begin{alertblock}{Critical Trap}
			CIA is Confidentiality, Integrity, Availability.
		\end{alertblock}
	\end{frame}
	
	\begin{frame}{Q62: Confidentiality Definition}
		\begin{block}{Question}
			Confidentiality ensures:
		\end{block}
		\begin{enumerate}[(A)]
			\item Data is always available
			\item Data has not been modified
			\item Data is only accessed by authorized users
			\item Data is backed up
		\end{enumerate}
	\end{frame}
	
	\begin{frame}{Q62: Explanation}
		\begin{block}{Correct Answer: C}
			\textbf{Confidentiality restricts access:}
			\begin{itemize}
				\item Only authorized users access data
				\item Enforced through access control and encryption
				\item Availability and integrity are separate
			\end{itemize}
		\end{block}
		\begin{alertblock}{Critical Trap}
			Confidentiality is about access, not availability.
		\end{alertblock}
	\end{frame}
	
	\begin{frame}{Q63: Integrity Definition}
		\begin{block}{Question}
			Integrity ensures:
		\end{block}
		\begin{enumerate}[(A)]
			\item Data is always available
			\item Data is accurate and complete
			\item Data is only accessed by authorized users
			\item Data is encrypted
		\end{enumerate}
	\end{frame}
	
	\begin{frame}{Q63: Explanation}
		\begin{block}{Correct Answer: B}
			\textbf{Integrity = accuracy and completeness:}
			\begin{itemize}
				\item Data is not altered improperly
				\item Hashing and controls verify integrity
				\item Different from confidentiality
			\end{itemize}
		\end{block}
		\begin{alertblock}{Critical Trap}
			Integrity is about correctness, not access.
		\end{alertblock}
	\end{frame}
	
	\begin{frame}{Q64: Availability Definition}
		\begin{block}{Question}
			Availability ensures:
		\end{block}
		\begin{enumerate}[(A)]
			\item Data is accurate
			\item Data is only accessed by authorized users
			\item Data is accessible when needed
			\item Data is encrypted
		\end{enumerate}
	\end{frame}
	
	\begin{frame}{Q64: Explanation}
		\begin{block}{Correct Answer: C}
			\textbf{Availability = access when needed:}
			\begin{itemize}
				\item Systems and data are up and reachable
				\item Redundancy, backups, DDoS protection support it
				\item Distinct from confidentiality and integrity
			\end{itemize}
		\end{block}
		\begin{alertblock}{Critical Trap}
			Availability is about uptime, not secrecy.
		\end{alertblock}
	\end{frame}
	
	\begin{frame}{Q65: Access Management Pillars}
		\begin{block}{Question}
			Which group defines the four pillars of access management?
		\end{block}
		\begin{enumerate}[(A)]
			\item CIA
			\item IAAA (Identification, Authentication, Authorization, Accountability)
			\item SWOT
			\item SMART
		\end{enumerate}
	\end{frame}
	
	\begin{frame}{Q65: Explanation}
		\begin{block}{Correct Answer: B}
			\textbf{IAAA = four pillars:}
			\begin{itemize}
				\item Identification = who you claim to be
				\item Authentication = proof of identity
				\item Authorization = what you can do
				\item Accountability = what you did
			\end{itemize}
		\end{block}
		\begin{alertblock}{Critical Trap}
			IAAA is the four-pillar model.
		\end{alertblock}
	\end{frame}
	
	\begin{frame}{Q66: MFA Definition}
		\begin{block}{Question}
			Multi-factor authentication requires:
		\end{block}
		\begin{enumerate}[(A)]
			\item Two passwords
			\item Two or more different authentication factors
			\item One password and a username
			\item Any two identifiers
		\end{enumerate}
	\end{frame}
	
	\begin{frame}{Q66: Explanation}
		\begin{block}{Correct Answer: B}
			\textbf{Different factors, not duplicates:}
			\begin{itemize}
				\item Factors are different types: know, have, are
				\item Two passwords are not MFA
				\item Username plus password is single-factor
			\end{itemize}
		\end{block}
		\begin{alertblock}{Critical Trap}
			Two of the same factor is not MFA.
		\end{alertblock}
	\end{frame}
	
	\begin{frame}{Q67: Least Privilege}
		\begin{block}{Question}
			What does the principle of least privilege require?
		\end{block}
		\begin{enumerate}[(A)]
			\item Give everyone admin rights
			\item Give users the minimum access needed to do their work
			\item Deny all access
			\item Rotate passwords
		\end{enumerate}
	\end{frame}
	
	\begin{frame}{Q67: Explanation}
		\begin{block}{Correct Answer: B}
			\textbf{Minimum access needed:}
			\begin{itemize}
				\item Minimum access reduces blast radius
				\item Admin rights increase risk
				\item Enforced via RBAC
			\end{itemize}
		\end{block}
		\begin{alertblock}{Critical Trap}
			Least privilege is minimum access needed.
		\end{alertblock}
	\end{frame}
	
	\begin{frame}{Q68: Segregation of Duties}
		\begin{block}{Question}
			Segregation of duties (SoD) exists to:
		\end{block}
		\begin{enumerate}[(A)]
			\item Speed up processes
			\item Prevent any single person from controlling an entire critical process
			\item Reduce hiring
			\item Avoid testing
		\end{enumerate}
	\end{frame}
	
	\begin{frame}{Q68: Explanation}
		\begin{block}{Correct Answer: B}
			\textbf{SoD distributes control:}
			\begin{itemize}
				\item Distributes control to prevent fraud and errors
				\item Common in finance and IT
				\item Reduces single-point abuse
			\end{itemize}
		\end{block}
		\begin{alertblock}{Critical Trap}
			SoD prevents one person from controlling everything.
		\end{alertblock}
	\end{frame}
	
	\begin{frame}{Q69: RBAC Definition}
		\begin{block}{Question}
			Role-based access control (RBAC) assigns permissions based on:
		\end{block}
		\begin{enumerate}[(A)]
			\item User requests
			\item Roles tied to job functions
			\item Departments only
			\item Password strength
		\end{enumerate}
	\end{frame}
	
	\begin{frame}{Q69: Explanation}
		\begin{block}{Correct Answer: B}
			\textbf{RBAC = roles, not individuals:}
			\begin{itemize}
				\item Roles map to job responsibilities
				\item Permissions assigned to roles
				\item Users inherit role permissions
			\end{itemize}
		\end{block}
		\begin{alertblock}{Critical Trap}
			RBAC is about roles, not individuals.
		\end{alertblock}
	\end{frame}
	
	\begin{frame}{Q70: Access Recertification}
		\begin{block}{Question}
			Why is periodic access recertification necessary?
		\end{block}
		\begin{enumerate}[(A)]
			\item To save money
			\item To verify that users still need their access rights
			\item To reduce passwords
			\item To satisfy vendors
		\end{enumerate}
	\end{frame}
	
	\begin{frame}{Q70: Explanation}
		\begin{block}{Correct Answer: B}
			\textbf{Unused access is risk:}
			\begin{itemize}
				\item Access accumulates over time
				\item Recertification removes unnecessary access
				\item Supports least privilege and compliance
			\end{itemize}
		\end{block}
		\begin{alertblock}{Critical Trap}
			Unused access is risk that must be recertified.
		\end{alertblock}
	\end{frame}
	
	\begin{frame}{Q71: Symmetric vs Asymmetric}
		\begin{block}{Question}
			Which statement BEST distinguishes symmetric and asymmetric encryption?
		\end{block}
		\begin{enumerate}[(A)]
			\item They are identical
			\item Symmetric uses one key; asymmetric uses a key pair
			\item Asymmetric uses one key; symmetric uses two
			\item Both use only one key
		\end{enumerate}
	\end{frame}
	
	\begin{frame}{Q71: Explanation}
		\begin{block}{Correct Answer: B}
			\textbf{Symmetric vs asymmetric:}
			\begin{itemize}
				\item Symmetric is fast, one shared key
				\item Asymmetric is slower, uses public/private keys
				\item Both are important
			\end{itemize}
		\end{block}
		\begin{alertblock}{Critical Trap}
			Symmetric is one key; asymmetric is a key pair.
		\end{alertblock}
	\end{frame}
	
	\begin{frame}{Q72: Hashing Purpose}
		\begin{block}{Question}
			What is hashing used for?
		\end{block}
		\begin{enumerate}[(A)]
			\item Encrypting large files for confidentiality
			\item Verifying integrity and storing passwords
			\item Sending emails
			\item Managing networks
		\end{enumerate}
	\end{frame}
	
	\begin{frame}{Q72: Explanation}
		\begin{block}{Correct Answer: B}
			\textbf{Hashing = integrity:}
			\begin{itemize}
				\item One-way hash verifies integrity
				\item Password storage uses hashes
				\item Hashing is not encryption
			\end{itemize}
		\end{block}
		\begin{alertblock}{Critical Trap}
			Hashing is not encryption; it is integrity.
		\end{alertblock}
	\end{frame}
	
	\begin{frame}{Q73: Digital Signature Provides}
		\begin{block}{Question}
			Which properties does a digital signature provide?
		\end{block}
		\begin{enumerate}[(A)]
			\item Confidentiality and availability
			\item Integrity, authenticity, and non-repudiation
			\item Anonymity
			\item Availability only
		\end{enumerate}
	\end{frame}
	
	\begin{frame}{Q73: Explanation}
		\begin{block}{Correct Answer: B}
			\textbf{Signature provides three properties:}
			\begin{itemize}
				\item Signed hash ensures integrity
				\item Signer identity provides authenticity
				\item Sender cannot deny
				\item Confidentiality needs encryption
			\end{itemize}
		\end{block}
		\begin{alertblock}{Critical Trap}
			Signatures do NOT provide confidentiality.
		\end{alertblock}
	\end{frame}
	
	\begin{frame}{Q74: Encryption vs Signature}
		\begin{block}{Question}
			To send a message with BOTH confidentiality AND non-repudiation, the sender should:
		\end{block}
		\begin{enumerate}[(A)]
			\item Only sign the message
			\item Only encrypt the message
			\item Sign AND encrypt the message
			\item Delete the message
		\end{enumerate}
	\end{frame}
	
	\begin{frame}{Q74: Explanation}
		\begin{block}{Correct Answer: C}
			\textbf{Sign and encrypt:}
			\begin{itemize}
				\item Encryption provides confidentiality
				\item Signature provides integrity and non-repudiation
				\item Both together provide all three
			\end{itemize}
		\end{block}
		\begin{alertblock}{Critical Trap}
			Confidentiality plus non-repudiation means sign and encrypt.
		\end{alertblock}
	\end{frame}
	
	\begin{frame}{Q75: PKI Trust Anchor}
		\begin{block}{Question}
			What is the trust anchor in a PKI?
		\end{block}
		\begin{enumerate}[(A)]
			\item The user's browser
			\item The Certificate Authority
			\item The DNS server
			\item The firewall
		\end{enumerate}
	\end{frame}
	
	\begin{frame}{Q75: Explanation}
		\begin{block}{Correct Answer: B}
			\textbf{The CA is the trust anchor:}
			\begin{itemize}
				\item CA issues and signs certificates
				\item Trust flows from the CA
				\item CA compromise breaks PKI trust
			\end{itemize}
		\end{block}
		\begin{alertblock}{Critical Trap}
			The CA is the trust anchor.
		\end{alertblock}
	\end{frame}
	
	\begin{frame}{Q76: CRL Purpose}
		\begin{block}{Question}
			What is the purpose of a Certificate Revocation List (CRL)?
		\end{block}
		\begin{enumerate}[(A)]
			\item Add new certificates
			\item List certificates that have been revoked before expiry
			\item Encrypt data
			\item Share public keys
		\end{enumerate}
	\end{frame}
	
	\begin{frame}{Q76: Explanation}
		\begin{block}{Correct Answer: B}
			\textbf{Revoked is not expired:}
			\begin{itemize}
				\item Revoked certs are no longer trusted
				\item CRLs and OCSP inform relying parties
				\item Expiry and revocation are different
			\end{itemize}
		\end{block}
		\begin{alertblock}{Critical Trap}
			Revoked is different from expired.
		\end{alertblock}
	\end{frame}
	
	\begin{frame}{Q77: Security Awareness Goal}
		\begin{block}{Question}
			What is the primary goal of security awareness training?
		\end{block}
		\begin{enumerate}[(A)]
			\item Show completion rates
			\item Build a security-conscious culture and reduce human risk
			\item Meet one audit requirement
			\item Replace technical controls
		\end{enumerate}
	\end{frame}
	
	\begin{frame}{Q77: Explanation}
		\begin{block}{Correct Answer: B}
			\textbf{Behavior change is the goal:}
			\begin{itemize}
				\item Behavior change is the goal
				\item Completion is a process metric
				\item Training complements controls
			\end{itemize}
		\end{block}
		\begin{alertblock}{Critical Trap}
			Awareness is about behavior, not completion.
		\end{alertblock}
	\end{frame}
	
	\begin{frame}{Q78: Phishing Defense}
		\begin{block}{Question}
			Which combination BEST defends against phishing?
		\end{block}
		\begin{enumerate}[(A)]
			\item Only firewalls
			\item Awareness, email filtering, MFA, and reporting channels
			\item Disabling email
			\item Blocking all links
		\end{enumerate}
	\end{frame}
	
	\begin{frame}{Q78: Explanation}
		\begin{block}{Correct Answer: B}
			\textbf{Phishing defense is layered:}
			\begin{itemize}
				\item Defense in depth reduces success
				\item MFA reduces credential theft impact
				\item Reporting catches missed attacks
			\end{itemize}
		\end{block}
		\begin{alertblock}{Critical Trap}
			Phishing defense is layered, not single.
		\end{alertblock}
	\end{frame}
	
	\begin{frame}{Q79: Insider Threat}
		\begin{block}{Question}
			Which control BEST addresses insider threats?
		\end{block}
		\begin{enumerate}[(A)]
			\item Only firewalls
			\item Least privilege, monitoring, and behavior analytics
			\item Disabling email
			\item Blocking all remote access
		\end{enumerate}
	\end{frame}
	
	\begin{frame}{Q79: Explanation}
		\begin{block}{Correct Answer: B}
			\textbf{Insider threat requires layered monitoring:}
			\begin{itemize}
				\item Least privilege limits damage
				\item Monitoring detects anomalies
				\item Behavior analytics catches unusual access
			\end{itemize}
		\end{block}
		\begin{alertblock}{Critical Trap}
			Insider threat requires layered monitoring.
		\end{alertblock}
	\end{frame}
	
	\begin{frame}{Q80: Data Privacy vs Security}
		\begin{block}{Question}
			Which statement BEST distinguishes data privacy from data security?
		\end{block}
		\begin{enumerate}[(A)]
			\item They are identical
			\item Privacy is about proper use and rights; security is about protecting data
			\item Privacy only applies to backups
			\item Security only applies to passwords
		\end{enumerate}
	\end{frame}
	
	\begin{frame}{Q80: Explanation}
		\begin{block}{Correct Answer: B}
			\textbf{Privacy is not the same as security:}
			\begin{itemize}
				\item Privacy is about rights and consent
				\item Security is about protection
				\item Security enables privacy
			\end{itemize}
		\end{block}
		\begin{alertblock}{Critical Trap}
			Privacy is not the same as security.
		\end{alertblock}
	\end{frame}
	
	\begin{frame}{Q81: GDPR Purpose}
		\begin{block}{Question}
			Which is the PRIMARY purpose of the EU GDPR?
		\end{block}
		\begin{enumerate}[(A)]
			\item Protect personal data and privacy of EU residents
			\item Reduce IT budgets
			\item Ban cloud computing
			\item Require encryption always
		\end{enumerate}
	\end{frame}
	
	\begin{frame}{Q81: Explanation}
		\begin{block}{Correct Answer: A}
			\textbf{GDPR protects EU data subjects:}
			\begin{itemize}
				\item Protects personal data and privacy
				\item Applies to any organization handling EU resident data
				\item Strong penalties for non-compliance
			\end{itemize}
		\end{block}
		\begin{alertblock}{Critical Trap}
			GDPR protects EU data subjects, not just EU companies.
		\end{alertblock}
	\end{frame}
	
	\begin{frame}{Q82: CCPA Beneficiary}
		\begin{block}{Question}
			Which group is a beneficiary of the CCPA?
		\end{block}
		\begin{enumerate}[(A)]
			\item US residents in general
			\item EU residents
			\item California residents
			\item Canadian residents
		\end{enumerate}
	\end{frame}
	
	\begin{frame}{Q82: Explanation}
		\begin{block}{Correct Answer: C}
			\textbf{CCPA is about California residents:}
			\begin{itemize}
				\item CCPA protects California residents
				\item GDPR covers EU residents
				\item Different laws, different regions
			\end{itemize}
		\end{block}
		\begin{alertblock}{Critical Trap}
			CCPA is California residents, not all US.
		\end{alertblock}
	\end{frame}
	
	\begin{frame}{Q83: HIPAA Breach Notification}
		\begin{block}{Question}
			Under HIPAA, when must a breach be reported to OCR?
		\end{block}
		\begin{enumerate}[(A)]
			\item More than 100 individuals
			\item Any IT incident
			\item More than 500 individuals
			\item More than 1,000 individuals
		\end{enumerate}
	\end{frame}
	
	\begin{frame}{Q83: Explanation}
		\begin{block}{Correct Answer: C}
			\textbf{HIPAA threshold is 500:}
			\begin{itemize}
				\item 500 or more individuals requires OCR reporting
				\item Fewer may still require reporting to individuals
				\item Threshold is a key HIPAA rule
			\end{itemize}
		\end{block}
		\begin{alertblock}{Critical Trap}
			HIPAA threshold is 500 for OCR notification.
		\end{alertblock}
	\end{frame}
	
	\begin{frame}{Q84: Data Subject Rights}
		\begin{block}{Question}
			Which is an example of a data subject right under GDPR?
		\end{block}
		\begin{enumerate}[(A)]
			\item Right to unlimited access
			\item Right to be forgotten
			\item Right to avoid all audits
			\item Right to skip training
		\end{enumerate}
	\end{frame}
	
	\begin{frame}{Q84: Explanation}
		\begin{block}{Correct Answer: B}
			\textbf{Right to be forgotten:}
			\begin{itemize}
				\item Right to erasure
				\item Right to access, correction, portability
				\item Enforced through processes
			\end{itemize}
		\end{block}
		\begin{alertblock}{Critical Trap}
			Right to be forgotten is a GDPR right.
		\end{alertblock}
	\end{frame}
	
	\begin{frame}{Q85: Privacy by Design}
		\begin{block}{Question}
			What does privacy by design require?
		\end{block}
		\begin{enumerate}[(A)]
			\item Privacy added after go-live
			\item Privacy embedded into system design from the start
			\item No data collected
			\item Ignoring regulations
		\end{enumerate}
	\end{frame}
	
	\begin{frame}{Q85: Explanation}
		\begin{block}{Correct Answer: B}
			\textbf{Privacy starts at design:}
			\begin{itemize}
				\item Privacy embedded early
				\item Cheaper and more effective
				\item Common in GDPR
			\end{itemize}
		\end{block}
		\begin{alertblock}{Critical Trap}
			Privacy by design starts at design, not after launch.
		\end{alertblock}
	\end{frame}
	
	\begin{frame}{Q86: Consent}
		\begin{block}{Question}
			Valid consent under GDPR must be:
		\end{block}
		\begin{enumerate}[(A)]
			\item Implied
			\item Freely given, specific, informed, and unambiguous
			\item Assumed
			\item Verbal only
		\end{enumerate}
	\end{frame}
	
	\begin{frame}{Q86: Explanation}
		\begin{block}{Correct Answer: B}
			\textbf{Consent must be clear:}
			\begin{itemize}
				\item Consent must be specific and clear
				\item Withdrawable at any time
				\item Recorded and demonstrable
			\end{itemize}
		\end{block}
		\begin{alertblock}{Critical Trap}
			GDPR consent is clear, explicit, and withdrawable.
		\end{alertblock}
	\end{frame}
	
	\begin{frame}{Q87: Data Breach Notification}
		\begin{block}{Question}
			What does GDPR require after a personal data breach?
		\end{block}
		\begin{enumerate}[(A)]
			\item Nothing
			\item Notify the supervisory authority within 72 hours if risk to data subjects
			\item Only notify marketing
			\item Delete all data
		\end{enumerate}
	\end{frame}
	
	\begin{frame}{Q87: Explanation}
		\begin{block}{Correct Answer: B}
			\textbf{72-hour rule:}
			\begin{itemize}
				\item 72 hours to notify supervisory authority
				\item Notify data subjects if high risk
				\item Document breaches
			\end{itemize}
		\end{block}
		\begin{alertblock}{Critical Trap}
			GDPR breach notification is 72 hours.
		\end{alertblock}
	\end{frame}
	
	\begin{frame}{Q88: Controller vs Processor}
		\begin{block}{Question}
			Which statement is TRUE about controllers and processors?
		\end{block}
		\begin{enumerate}[(A)]
			\item They are the same
			\item Controllers decide purposes; processors act on instructions
			\item Processors decide purposes
			\item Controllers only act on instructions
		\end{enumerate}
	\end{frame}
	
	\begin{frame}{Q88: Explanation}
		\begin{block}{Correct Answer: B}
			\textbf{Controller decides; processor executes:}
			\begin{itemize}
				\item Controller decides purpose and means
				\item Processor handles data on behalf
				\item Different legal obligations
			\end{itemize}
		\end{block}
		\begin{alertblock}{Critical Trap}
			Controller decides; processor executes.
		\end{alertblock}
	\end{frame}
	
	\begin{frame}{Q89: Encryption at Rest}
		\begin{block}{Question}
			What is the primary benefit of encryption at rest?
		\end{block}
		\begin{enumerate}[(A)]
			\item Faster performance
			\item Protecting data if storage media is lost or stolen
			\item Reducing backup costs
			\item Simplifying access
		\end{enumerate}
	\end{frame}
	
	\begin{frame}{Q89: Explanation}
		\begin{block}{Correct Answer: B}
			\textbf{At-rest encryption protects lost media:}
			\begin{itemize}
				\item Protects stored data
				\item Reduces risk of exposure
				\item Complementary to encryption in transit
			\end{itemize}
		\end{block}
		\begin{alertblock}{Critical Trap}
			At-rest encryption protects lost media, not live misuse.
		\end{alertblock}
	\end{frame}
	
	\begin{frame}{Q90: Encryption in Transit}
		\begin{block}{Question}
			What does encryption in transit protect against?
		\end{block}
		\begin{enumerate}[(A)]
			\item Lost laptops
			\item Eavesdropping and interception on networks
			\item Wrongly configured firewalls
			\item Lost passwords
		\end{enumerate}
	\end{frame}
	
	\begin{frame}{Q90: Explanation}
		\begin{block}{Correct Answer: B}
			\textbf{In-transit encryption protects network paths:}
			\begin{itemize}
				\item Protects data moving between systems
				\item Uses TLS, VPNs, IPSec
				\item Complementary to at-rest encryption
			\end{itemize}
		\end{block}
		\begin{alertblock}{Critical Trap}
			In-transit encryption protects network paths.
		\end{alertblock}
	\end{frame}
	
	\begin{frame}{Q91: Key Management}
		\begin{block}{Question}
			Why is key management critical in cryptography?
		\end{block}
		\begin{enumerate}[(A)]
			\item Keys are cheap
			\item Weak key management defeats the strongest encryption
			\item Keys are only for backups
			\item Keys are optional
		\end{enumerate}
	\end{frame}
	
	\begin{frame}{Q91: Explanation}
		\begin{block}{Correct Answer: B}
			\textbf{Strong encryption plus weak key management equals broken security:}
			\begin{itemize}
				\item Keys protect encrypted data
				\item Leaked keys break encryption
				\item Rotation, storage, and access control are essential
			\end{itemize}
		\end{block}
		\begin{alertblock}{Critical Trap}
			Strong encryption plus weak key management equals broken security.
		\end{alertblock}
	\end{frame}
	
	\begin{frame}{Q92: HSM Purpose}
		\begin{block}{Question}
			What is the PRIMARY purpose of a Hardware Security Module?
		\end{block}
		\begin{enumerate}[(A)]
			\item Speed up networks
			\item Securely store and manage cryptographic keys
			\item Back up files
			\item Manage DNS
		\end{enumerate}
	\end{frame}
	
	\begin{frame}{Q92: Explanation}
		\begin{block}{Correct Answer: B}
			\textbf{HSM protects keys:}
			\begin{itemize}
				\item HSMs protect keys
				\item Cryptographic operations occur inside
				\item Keys do not leave the HSM
			\end{itemize}
		\end{block}
		\begin{alertblock}{Critical Trap}
			HSMs protect keys, not networks.
		\end{alertblock}
	\end{frame}
	
	\begin{frame}{Q93: Zero Trust Principle}
		\begin{block}{Question}
			What does the Zero Trust security model assume?
		\end{block}
		\begin{enumerate}[(A)]
			\item Trust everything inside the network
			\item Trust nothing by default and verify every access
			\item Trust all employees
			\item Trust all vendors
		\end{enumerate}
	\end{frame}
	
	\begin{frame}{Q93: Explanation}
		\begin{block}{Correct Answer: B}
			\textbf{Verify, do not trust by default:}
			\begin{itemize}
				\item No implicit trust based on location
				\item Continuous verification
				\item Least privilege and segmentation
			\end{itemize}
		\end{block}
		\begin{alertblock}{Critical Trap}
			Zero Trust means verify, not trust by default.
		\end{alertblock}
	\end{frame}
	
	\begin{frame}{Q94: Defense in Depth}
		\begin{block}{Question}
			What does defense in depth provide?
		\end{block}
		\begin{enumerate}[(A)]
			\item A single strong control
			\item Multiple layered controls reducing single-point failure
			\item More vendors
			\item Less budget
		\end{enumerate}
	\end{frame}
	
	\begin{frame}{Q94: Explanation}
		\begin{block}{Correct Answer: B}
			\textbf{Multiple layers, not one:}
			\begin{itemize}
				\item Layers protect each other
				\item Independent controls reduce risk of failure
				\item Attackers must defeat multiple defenses
			\end{itemize}
		\end{block}
		\begin{alertblock}{Critical Trap}
			Defense in depth is multiple layers, not one.
		\end{alertblock}
	\end{frame}
	
	\begin{frame}{Q95: Security Control Types}
		\begin{block}{Question}
			Which grouping BEST classifies security controls?
		\end{block}
		\begin{enumerate}[(A)]
			\item Preventive, detective, corrective
			\item Red, green, blue
			\item Fast, slow, medium
			\item High, medium, low
		\end{enumerate}
	\end{frame}
	
	\begin{frame}{Q95: Explanation}
		\begin{block}{Correct Answer: A}
			\textbf{Preventive, detective, corrective:}
			\begin{itemize}
				\item Preventive = stop events
				\item Detective = identify events
				\item Corrective = reduce impact
			\end{itemize}
		\end{block}
		\begin{alertblock}{Critical Trap}
			Preventive, detective, corrective are core classes.
		\end{alertblock}
	\end{frame}
	
	\begin{frame}{Q96: Vulnerability vs Threat vs Risk}
		\begin{block}{Question}
			Which statement BEST distinguishes vulnerability, threat, and risk?
		\end{block}
		\begin{enumerate}[(A)]
			\item They are identical
			\item Vulnerability is weakness; threat is potential cause; risk is likelihood and impact
			\item Threat is weakness; risk is cause; vulnerability is impact
			\item Risk is only about costs
		\end{enumerate}
	\end{frame}
	
	\begin{frame}{Q96: Explanation}
		\begin{block}{Correct Answer: B}
			\textbf{Vulnerability, threat, and risk are different:}
			\begin{itemize}
				\item Vulnerability = weakness
				\item Threat = potential attacker or event
				\item Risk = likelihood and impact
			\end{itemize}
		\end{block}
		\begin{alertblock}{Critical Trap}
			Vulnerability, threat, and risk are different.
		\end{alertblock}
	\end{frame}
	
	\begin{frame}{Q97: Patch Management}
		\begin{block}{Question}
			What is the primary goal of patch management?
		\end{block}
		\begin{enumerate}[(A)]
			\item Reduce costs
			\item Remediate known vulnerabilities in a timely, controlled manner
			\item Slow down systems
			\item Add features
		\end{enumerate}
	\end{frame}
	
	\begin{frame}{Q97: Explanation}
		\begin{block}{Correct Answer: B}
			\textbf{Patch management is about vulnerabilities:}
			\begin{itemize}
				\item Address known vulnerabilities
				\item Controlled rollout reduces disruption
				\item Unpatched systems are prime targets
			\end{itemize}
		\end{block}
		\begin{alertblock}{Critical Trap}
			Patch management is about vulnerabilities, not features.
		\end{alertblock}
	\end{frame}
	
	\begin{frame}{Q98: Vulnerability Scoring}
		\begin{block}{Question}
			What is CVSS used for?
		\end{block}
		\begin{enumerate}[(A)]
			\item Scoring the severity of vulnerabilities
			\item Scoring employees
			\item Measuring backups
			\item Calculating budgets
		\end{enumerate}
	\end{frame}
	
	\begin{frame}{Q98: Explanation}
		\begin{block}{Correct Answer: A}
			\textbf{CVSS scores vulnerabilities:}
			\begin{itemize}
				\item Common Vulnerability Scoring System
				\item Quantifies severity
				\item Supports prioritization
			\end{itemize}
		\end{block}
		\begin{alertblock}{Critical Trap}
			CVSS scores vulnerabilities, not risks alone.
		\end{alertblock}
	\end{frame}
	
	\begin{frame}{Q99: Incident Response Stages}
		\begin{block}{Question}
			Which is a common stage in incident response?
		\end{block}
		\begin{enumerate}[(A)]
			\item Denial
			\item Preparation, detection, containment, eradication, recovery, lessons learned
			\item Marketing
			\item Budgeting
		\end{enumerate}
	\end{frame}
	
	\begin{frame}{Q99: Explanation}
		\begin{block}{Correct Answer: B}
			\textbf{IR is a cycle:}
			\begin{itemize}
				\item Structured stages ensure response quality
				\item Detection and containment are critical
				\item Lessons learned improve future response
			\end{itemize}
		\end{block}
		\begin{alertblock}{Critical Trap}
			Incident response is a cycle, not a single step.
		\end{alertblock}
	\end{frame}
	
	\begin{frame}{Q100: Ultimate Security Principle}
		\begin{block}{Question}
			What is the ultimate principle of security governance?
		\end{block}
		\begin{enumerate}[(A)]
			\item Buy the newest tools
			\item Align security with business objectives and manage risk continuously
			\item Maximize controls
			\item Eliminate all risk
		\end{enumerate}
	\end{frame}
	
	\begin{frame}{Q100: Explanation}
		\begin{block}{Correct Answer: B}
			\textbf{Security exists to enable the business:}
			\begin{itemize}
				\item Security supports business goals
				\item Risk is managed continuously
				\item Perfection is impossible; optimization is the goal
			\end{itemize}
		\end{block}
		\begin{alertblock}{Critical Trap}
			Security exists to enable the business, not block it.
		\end{alertblock}
	\end{frame}
	
	\begin{frame}{Exam Mindset Summary}
		\begin{block}{Remember}
			\begin{itemize}
				\item \textbf{CIA triad:} Confidentiality, Integrity, Availability
				\item \textbf{IAAA:} Identification, Authentication, Authorization, Accountability
				\item \textbf{Encryption vs Hashing:} Encryption is reversible; hashing is not
				\item \textbf{Signatures} provide integrity, authenticity, non-repudiation -- NOT confidentiality
				\item \textbf{Confidentiality + Non-repudiation} = sign AND encrypt
				\item \textbf{Privacy} is about rights; \textbf{security} is about protection
				\item \textbf{Zero Trust} = verify, do not trust by default
				\item \textbf{Defense in depth} = layered controls
				\item \textbf{Patch management} addresses known vulnerabilities
				\item \textbf{Quantum} threatens current public-key encryption
			\end{itemize}
		\end{block}
	\end{frame}
	
	\begin{frame}{Thank You}
		\begin{center}
			{\Huge \textbf{Good Luck on Your CRISC Exam!}}
		\end{center}
	\end{frame}
	
\end{document}